\section{Limits and Responsible Use}
\begin{frame}{}
    \LARGE Part 7: \textbf{Limits and Responsible Use}
\end{frame}

\begin{frame}{AI Makes Mistakes --- Confidently}
    \begin{itemize}
        \setlength\itemsep{0.8em}
        \item A language model is trained to produce text that \textbf{sounds right}. That is not the same as text that \textbf{is right}.
        \item It can invent facts, names, numbers and references. This is called a \textbf{hallucination}.
        \item It uses the same confident tone whether it is right or wrong.
        \item Its knowledge stops at the date its training data was collected.
    \end{itemize}

    \vspace{1em}
    \caveat{\textbf{Rule of thumb:} the more the answer matters, the more you must check it
    against a source you trust.}
\end{frame}

\begin{frame}{AI Learns Our Biases}
    \textit{A model learns from past data. If the past was unfair, the model repeats it.}

    \vspace{0.6em}
    \begin{columns}[T]
        \begin{column}{0.48\textwidth}
            \begin{block}{A hiring tool}
                \small
                A company trained a system on ten years of its own past hiring, mostly of men.
                According to Reuters, it learned to mark down CVs containing the word ``women's''.
                The project was abandoned.
            \end{block}
        \end{column}
        \begin{column}{0.48\textwidth}
            \begin{block}{Face analysis}
                \small
                A 2018 study tested three commercial systems that guess gender from a face.
                Errors were under 1\% for lighter-skinned men, and up to 35\% for
                darker-skinned women.
            \end{block}
        \end{column}
    \end{columns}

    \vspace{0.6em}
    \centering
    \textit{The system did not ``decide'' to be unfair. It copied the patterns in its data.}

    \vspace{0.4em}
    {\scriptsize J.~Dastin, \emph{Reuters}, 10 October 2018. \quad
    J.~Buolamwini and T.~Gebru, ``Gender Shades,'' \emph{PMLR} 81:77--91, 2018.}
\end{frame}

\begin{frame}{Myths and Reality}
    \centering
    \renewcommand{\arraystretch}{1.6}
    \begin{tabular}{@{} >{\raggedright\arraybackslash}p{0.38\textwidth}
                        >{\raggedright\arraybackslash}p{0.52\textwidth} @{}}
        \toprule
        \textbf{Myth} & \textbf{Reality} \\
        \midrule
        ``AI thinks like a person.'' & It finds patterns in data. It has no feelings or intentions. \\
        ``AI is always objective.'' & It reflects the data it was trained on, including its flaws. \\
        ``AI is always right.'' & It makes mistakes, and often sounds sure when it does. \\
        ``AI will take every job.'' & It changes tasks within jobs. People who use it well gain an advantage. \\
        ``You must be a maths expert to use AI.'' & Using AI tools well is mostly about clear thinking and checking results. \\
        \bottomrule
    \end{tabular}
\end{frame}

\begin{frame}{Using AI Responsibly}
    \begin{enumerate}
        \setlength\itemsep{0.9em}
        \item \textbf{Check important answers.} Treat the output as a draft, not as the truth.
        \item \textbf{Protect private information.} Do not paste passwords, personal data or confidential documents into public AI tools.
        \item \textbf{Be honest about using it.} Follow the rules of your school or workplace.
        \item \textbf{Stay responsible.} If you send it, sign it or submit it, it is yours --- not the AI's.
        \item \textbf{Ask who could be affected.} Would this be fair to everyone it touches?
    \end{enumerate}
\end{frame}
